\documentclass{amsart}
\usepackage[T1]{fontenc}
\usepackage{amsmath,amssymb,amsthm,hyperref}
\usepackage{algorithm}
\usepackage{algpseudocode}
\usepackage{bm}

\newtheorem{theorem}{Theorem}
\newtheorem{lemma}{Lemma}
\newtheorem{proposition}{Proposition}
\newtheorem{corollary}{Corollary}
\newtheorem{definition}{Definition}
\newtheorem{remark}{Remark}

\newcommand{\Sym}{\mathbb{S}}
\newcommand{\R}{\mathbb{R}}
\newcommand{\tr}{\operatorname{tr}}
\newcommand{\diag}{\operatorname{diag}}
\newcommand{\norm}[1]{\left\lVert#1\right\rVert}

\begin{document}

\title{A structure-exploiting interior-point algorithm for periodic linear matrix inequalities}
\author{solver-agent}
\date{}
\maketitle

\begin{abstract}
We give a complete algorithm that decides feasibility of a periodic linear matrix
inequality (periodic LMI) system, produces a $P$-periodic symmetric solution when
one exists, and never forms the lifted $nP\times nP$ block-diagonal LMI. The
algorithm is a barrier (interior-point) path-following method specialised to the
periodic structure. Its central component is a structural lemma showing that the
Newton system at every iteration is a \emph{symmetric positive-definite
block-cyclic-tridiagonal} system, which we solve exactly in $O(P)$ block operations
by a non-cyclic block factorisation corrected with a rank-$2m$
Sherman--Morrison--Woodbury update. We prove global convergence with the standard
self-concordant iteration bound $O(\sqrt{nP}\,\log(1/\varepsilon))$ outer Newton
steps, and we establish a per-iteration cost of $O\!\big(P(m^2 n^2 + m\,n^3 + n^3)\big)$
flops and $O(P(m + n^2))$ storage. We describe a periodic-Schur/periodic-QZ
based realisation of the block solves that is structure preserving and backward
stable.
\end{abstract}

\section{Formalisation}\label{sec:form}

Throughout, $\Sym^{n}$ denotes the real symmetric $n\times n$ matrices, $A\succ0$
(resp.\ $A\succeq0$) means $A$ is symmetric positive (semi)definite, and indices in
$k$ are read modulo $P$ (so $X(P)\equiv X(0)$, $A(k+P)\equiv A(k)$, etc.).

We are given $P$-periodic data: for each $k\in\{0,\dots,P-1\}$ a finite list of
symmetric matrices
\[
  F_0(k),F_1(k),\dots,F_m(k)\in\Sym^{n},
\]
and scalar decision variables collected into the vector
$x(k)=(x_1(k),\dots,x_m(k))^{\!\top}\in\R^{m}$, with the periodicity convention
$x(k+P)=x(k)$. Define the affine matrix map
\begin{equation}\label{eq:Gmap}
  G(k,x(k)) \;:=\; F_0(k) + \sum_{i=1}^{m} x_i(k)\,F_i(k)\ \in\ \Sym^{n}.
\end{equation}

We will also allow \emph{dynamic coupling} between consecutive instants $k$ and
$k+1$, which is the feature that distinguishes a genuine periodic LMI from $P$
independent LMIs. In its most general form the coupling is described by
$P$-periodic data and the coupled constraint
\begin{equation}\label{eq:coupled}
  G(k,x(k)) \;+\; C\big(k,\,x(k),\,x(k+1)\big)\ \succ\ 0,
  \qquad k=0,\dots,P-1,
\end{equation}
where $C$ is \emph{jointly affine} in $\big(x(k),x(k+1)\big)$ and is the only
mechanism through which different time instants interact. The canonical examples
are:

\begin{itemize}
\item \textbf{Periodic Lyapunov inequality.}
With matrix variable $X(k)=X(k)^{\!\top}$, $X(k+P)=X(k)$,
\begin{equation}\label{eq:lyap}
  -\big(A(k)^{\!\top}X(k+1)A(k)-X(k)+Q(k)\big)\succ0,\qquad X(k)\succ0,
\end{equation}
which is of the form \eqref{eq:coupled} with $x(k)$ a vectorisation of $X(k)$, the
term $-X(k)+Q(k)$ entering through $G$ and the coupling term
$C(k,x(k),x(k+1))=-A(k)^{\!\top}X(k+1)A(k)$ entering through $C$.
\item \textbf{Bounded-real / KYP variants}, obtained by appending the channel
matrices $B(k),C(k),D(k)$ to the Lyapunov dissipation inequality; these are again of
the form \eqref{eq:coupled} with $C$ depending only on $\big(X(k),X(k+1)\big)$.
\end{itemize}

We treat \eqref{eq:coupled} because it strictly contains the uncoupled system
$G(k,x(k))\succ0$ (set $C\equiv0$). The unifying abstraction is the following.

\begin{definition}[Periodic LMI feasibility problem, PLMI]\label{def:plmi}
Let $\mathcal{X}:=\R^{mP}$ with a block variable $\bm x=(x(0),\dots,x(P-1))$. For
each $k$ let $\mathcal{A}_k:\mathcal X\to\Sym^{n}$ be the affine map
\begin{equation}\label{eq:Ak}
  \mathcal{A}_k(\bm x)\;:=\;G(k,x(k))+C(k,x(k),x(k+1)),
\end{equation}
which by construction depends only on the two blocks $x(k)$ and $x(k+1)$. The PLMI
asks to decide whether the open convex set
\begin{equation}\label{eq:feasset}
  \mathcal F\;:=\;\Big\{\bm x\in\mathcal X:\ \mathcal{A}_k(\bm x)\succ0\ \text{for all }k=0,\dots,P-1\Big\}
\end{equation}
is nonempty and, if so, to return a point $\bm x\in\mathcal F$ (equivalently the
matrices $X(0),\dots,X(P-1)$ in the Lyapunov instance).
\end{definition}

The naive ``lifted'' reformulation stacks all constraints into one block-diagonal
LMI $\diag_k\mathcal A_k(\bm x)\succ0$ of size $nP\times nP$ and applies a general
SDP solver, costing $O\!\big((nP)^3\big)$ or worse per iteration and
$O\!\big((nP)^2\big)$ storage. Our goal is to do strictly better in $P$.

\medskip
\noindent\textbf{Standing nondegeneracy assumption (A).}
\emph{The feasible set $\mathcal F$, if nonempty, has nonempty interior and is
bounded after the homogenisation of \S\ref{sec:phaseI}.} This is the usual
Slater/boundedness hypothesis under which interior-point feasibility solvers are
formulated; it can always be enforced by adding a bounding constraint $\norm{\bm
x}\le R$ (which preserves all structural properties below, being block-diagonal in
$k$). We make it explicit so that the convergence statements are unconditional.

\section{The barrier and its derivatives}\label{sec:barrier}

For a convex feasibility problem the cleanest globally convergent method is the
\emph{analytic-center / barrier path-following} method built on a self-concordant
barrier. We recall the two facts we use; both are classical
(Nesterov--Nemirovski, \emph{Interior-Point Polynomial Algorithms in Convex
Programming}, SIAM 1994, Thm.~2.1.1 and Prop.~5.1.4).

\begin{lemma}[Log-det barrier]\label{lem:sc}
The function $\phi(S)=-\log\det S$ on the cone $\Sym^{n}_{++}$ is a self-concordant
barrier with parameter $\nu=n$: it satisfies, for all $S\succ0$ and all
$H\in\Sym^n$,
\[
  |D^3\phi(S)[H,H,H]|\le 2\big(D^2\phi(S)[H,H]\big)^{3/2},\qquad
  \big(D\phi(S)[H]\big)^2\le n\, D^2\phi(S)[H,H].
\]
If $\mathcal A:\R^{N}\to\Sym^{n}$ is affine, then $\bm x\mapsto\phi(\mathcal A(\bm
x))$ is self-concordant on $\{\bm x:\mathcal A(\bm x)\succ0\}$, with the same barrier
parameter $n$, and a finite sum of self-concordant barriers is self-concordant with
parameter equal to the sum of the parameters.
\end{lemma}

We attach one log-det term per constraint. Write the residual
$S_k(\bm x):=\mathcal A_k(\bm x)\succ0$ and define the barrier
\begin{equation}\label{eq:barrier}
  F(\bm x)\;:=\;\sum_{k=0}^{P-1} \phi\big(S_k(\bm x)\big)
            \;=\;-\sum_{k=0}^{P-1}\log\det\mathcal A_k(\bm x).
\end{equation}
In the Lyapunov instance \eqref{eq:lyap}, where the extra constraint $X(k)\succ0$ is
present, one simply adds the further log-det terms $-\sum_k\log\det X(k)$; this only
adds block-\emph{diagonal} (in $k$) curvature and does not affect any structural
claim below, so we omit it from the notation. By Lemma~\ref{lem:sc} the total
barrier parameter is
\begin{equation}\label{eq:nu}
  \nu\;\le\;nP\qquad(\text{or }2nP\text{ in the Lyapunov instance with }X(k)\succ0).
\end{equation}

\paragraph{Gradient and Hessian in the periodic structure.}
Using $D\phi(S)[H]=-\tr(S^{-1}H)$ and
$D^2\phi(S)[H,H']=\tr(S^{-1}HS^{-1}H')$ (standard log-det derivatives), and writing
$\partial_{x_i(k)}S_\ell$ for the (constant) derivative matrix of the residual
$S_\ell$ with respect to the scalar $x_i(k)$, we obtain
\begin{align}
  \big[\nabla F\big]_{i,k}
   &= -\sum_{\ell=0}^{P-1}\tr\!\big(S_\ell^{-1}\,\partial_{x_i(k)}S_\ell\big),
   \label{eq:grad}\\
  \big[\nabla^2 F\big]_{(i,k),(j,k')}
   &= \sum_{\ell=0}^{P-1}
      \tr\!\big(S_\ell^{-1}\,\partial_{x_i(k)}S_\ell\;
                S_\ell^{-1}\,\partial_{x_j(k')}S_\ell\big).
   \label{eq:hess}
\end{align}
The locality of \eqref{eq:coupled}--\eqref{eq:Ak} now produces the key sparsity.

\begin{lemma}[Block-cyclic-tridiagonal Hessian]\label{lem:struct}
Group the variables by time index, $\bm x=(x(0),\dots,x(P-1))$ with blocks
$x(k)\in\R^{m}$. Then $S_\ell$ depends only on $x(\ell)$ and $x(\ell+1)$. Hence in
\eqref{eq:hess} a term with summation index $\ell$ is nonzero only if both $(i,k)$
and $(j,k')$ index variables among $\{x(\ell),x(\ell+1)\}$. Consequently the
$m\times m$ block $\big[\nabla^2F\big]_{k,k'}$ is nonzero only for
$k'\in\{k-1,k,k+1\}\pmod P$. That is, $\nabla^2 F$ is a symmetric, block-cyclic
(periodic) block-tridiagonal matrix with block size $m$,
\begin{equation}\label{eq:HBT}
  \nabla^2F=
  \begin{pmatrix}
    D_0 & E_0 &        &        & E_{P-1}^{\!\top}\\
    E_0^{\!\top} & D_1 & E_1 &        & \\
        & E_1^{\!\top} & D_2 & \ddots & \\
        &        & \ddots & \ddots & E_{P-2}\\
    E_{P-1} &     &        & E_{P-2}^{\!\top} & D_{P-1}
  \end{pmatrix},
  \quad D_k,E_k\in\R^{m\times m}.
\end{equation}
Moreover $\nabla^2 F\succ0$ on the feasible set: it is the Hessian of a strictly
convex self-concordant barrier whose domain is the nonempty open convex set
$\mathcal F$ (by Assumption~(A) and Lemma~\ref{lem:sc}), so it is positive definite
at every $\bm x\in\mathcal F$.
\end{lemma}

\begin{proof}
The dependency claim is immediate from \eqref{eq:Ak}: $S_\ell=\mathcal A_\ell(\bm
x)=G(\ell,x(\ell))+C(\ell,x(\ell),x(\ell+1))$ is a function of $x(\ell),x(\ell+1)$
only, so $\partial_{x_i(k)}S_\ell\neq0$ forces $k\in\{\ell,\ell+1\}$. Therefore the
summand in \eqref{eq:hess} indexed by $\ell$ vanishes unless
$\{k,k'\}\subseteq\{\ell,\ell+1\}$, which forces $|k-k'|\le1$ modulo $P$. The block
$(k,k+1)$ (and its transpose $(k+1,k)$) receives contributions only from $\ell=k$;
the diagonal block $(k,k)$ from $\ell=k$ and $\ell=k-1$; the corner block $(P-1,0)$
from $\ell=P-1$. This is exactly the pattern \eqref{eq:HBT}.

For positive definiteness, fix $\bm x\in\mathcal F$ and a direction $\bm h\neq0$.
By \eqref{eq:hess},
\[
  \bm h^{\!\top}\nabla^2 F(\bm x)\,\bm h
  =\sum_{\ell=0}^{P-1}\tr\!\big(S_\ell^{-1}\,(\textstyle\sum_{i,k}h_{i,k}\partial_{x_i(k)}S_\ell)\,
                                 S_\ell^{-1}\,(\textstyle\sum_{j,k'}h_{j,k'}\partial_{x_j(k')}S_\ell)\big)
  =\sum_{\ell=0}^{P-1}\norm{S_\ell^{-1/2}\,\dot S_\ell\,S_\ell^{-1/2}}_F^2\ge0,
\]
where $\dot S_\ell:=\sum_{i,k}h_{i,k}\partial_{x_i(k)}S_\ell$ is the directional
derivative of $S_\ell$ and $\norm{\cdot}_F$ is the Frobenius norm; each summand is a
squared norm because $\tr(S^{-1}\dot S S^{-1}\dot S)=\norm{S^{-1/2}\dot S
S^{-1/2}}_F^2$ for symmetric $\dot S$. The quadratic form vanishes iff $\dot
S_\ell=0$ for all $\ell$, i.e.\ iff $\bm h$ lies in the kernel of the full affine map
$\bm x\mapsto(S_0,\dots,S_{P-1})$. Under Assumption~(A) this map is injective on
$\mathcal X$ (a nonzero kernel direction would make $\mathcal F$ unbounded along a
line, contradicting boundedness after homogenisation), hence $\dot S_\ell=0\ \forall
\ell\Rightarrow\bm h=0$. Therefore $\bm h^{\!\top}\nabla^2F(\bm x)\bm h>0$ for all
$\bm h\neq0$, i.e.\ $\nabla^2F(\bm x)\succ0$.
\end{proof}

\begin{remark}
The lifted formulation would treat $\nabla^2F$ as a dense $mP\times mP$ matrix and
factor it in $O\!\big((mP)^3\big)$. Lemma~\ref{lem:struct} shows it is in fact
banded plus rank-$2m$ cyclic correction, which is the structure we exploit in
\S\ref{sec:linalg}.
\end{remark}

\section{The structured Newton solver}\label{sec:linalg}

The single linear-algebra primitive used at every Newton step is: given the SPD
block-cyclic-tridiagonal matrix $H=\nabla^2F$ of \eqref{eq:HBT} and a right-hand
side $g\in\R^{mP}$, solve $H\,\Delta=g$. We do this in $O(P)$ block operations,
never assembling the $mP\times mP$ matrix.

\paragraph{Reduction to non-cyclic plus low-rank.}
Split $H=T+UKU^{\!\top}$, where $T$ is the \emph{non-cyclic} block-tridiagonal matrix
obtained from $H$ by deleting the two corner blocks $E_{P-1}$ (at block position
$(P-1,0)$) and $E_{P-1}^{\!\top}$ (at $(0,P-1)$), and the correction reinstates them
through
\begin{equation}\label{eq:UKV}
  U=\big[\,e_{0}\otimes I_m\ \ e_{P-1}\otimes I_m\,\big]\in\R^{mP\times2m},
  \qquad
  K=\begin{pmatrix}0 & E_{P-1}^{\!\top}\\ E_{P-1} & 0\end{pmatrix}\in\R^{2m\times2m},
\end{equation}
so that $UKU^{\!\top}$ has the block $E_{P-1}^{\!\top}$ at $(0,P-1)$ and $E_{P-1}$ at
$(P-1,0)$ and is zero elsewhere; here $e_k\in\R^{P}$ is the $k$-th unit vector and
$\otimes$ the Kronecker product. (When $E_{P-1}=0$, e.g.\ in genuinely non-cyclic
relaxations, $T=H$ and no correction is needed.) Note $\widetilde D_k=D_k$ for
$k\neq 0,P-1$ and $\widetilde D_0=D_0$, $\widetilde D_{P-1}=D_{P-1}$, since the
corner deletion only removes off-diagonal blocks; thus $T$ has the same diagonal
blocks as $H$ and is itself the Hessian of the strictly convex barrier on the
\emph{open chain} relaxation, hence SPD on $\mathcal F$ provided the chain
relaxation has the same injectivity; in any case $T\succ0$ whenever $H\succ0$ and
$UKU^{\!\top}\preceq0$ is avoided by the standard regularisation below.

\paragraph{Block tridiagonal solve (block Thomas / LDL).}
For the non-cyclic $T$ with diagonal blocks $\widetilde D_k$ and upper blocks $E_k$,
$k=0,\dots,P-2$, a single forward/backward block sweep computes the solution. Define
the block factorisation sweep
\begin{equation}\label{eq:thomas}
  \begin{aligned}
  &M_0=\widetilde D_0,\qquad
   M_k=\widetilde D_k - E_{k-1}^{\!\top}M_{k-1}^{-1}E_{k-1}\ \ (k\ge1),\\
  &y_0=g_0,\qquad y_k=g_k-E_{k-1}^{\!\top}M_{k-1}^{-1}y_{k-1}\ \ (k\ge1),\\
  &\Delta_{P-1}=M_{P-1}^{-1}y_{P-1},\qquad
   \Delta_k=M_k^{-1}\big(y_k-E_k\,\Delta_{k+1}\big)\ (k\le P-2).
  \end{aligned}
\end{equation}
The forward sweep is precisely the block LDL$^{\!\top}$ factorisation of $T$: its
pivots $M_k$ are the successive Schur complements of the leading principal block
submatrices of $T$. If $T\succ0$ these Schur complements are all positive definite,
hence invertible, so the sweep is well defined; in our implementation each $M_k$ is
factored by a pivoted (Bunch--Kaufman) symmetric factorisation, which is backward
stable (\S\ref{sec:schur}). If $T$ fails to be positive definite we replace it by
the SPD regularisation $T_\rho:=T+\rho U U^{\!\top}$ with $\rho$ chosen so that
$T_\rho\succ0$ (always possible since $H=T_\rho+(K-\rho I)$-type corrections remain a
finite Woodbury update); this preserves all complexity bounds and keeps the final
result exact, as shown next.

\paragraph{Cyclic correction by Sherman--Morrison--Woodbury.}
With the factorisation of $T$ (equivalently $T_\rho$) computed once, $T^{-1}$ acts as
an operator: each application is one back/forward block sweep of \eqref{eq:thomas}.
The Woodbury identity gives
\begin{equation}\label{eq:woodbury}
  H^{-1}g
   = T^{-1}g - T^{-1}U\big(K^{-1}+U^{\!\top}T^{-1}U\big)^{-1}U^{\!\top}T^{-1}g ,
\end{equation}
and when $K$ is singular one uses the algebraically equivalent form
\begin{equation}\label{eq:woodbury2}
  H^{-1}g
   = T^{-1}g - T^{-1}U\big(I_{2m}+KU^{\!\top}T^{-1}U\big)^{-1}KU^{\!\top}T^{-1}g ,
\end{equation}
valid whenever $H$ and $T$ are invertible. Here $U^{\!\top}T^{-1}U\in\R^{2m\times2m}$,
so the inner solve is a tiny $2m\times2m$ system.

\begin{lemma}[Exact structured solve]\label{lem:exact}
Formulas \eqref{eq:thomas}--\eqref{eq:woodbury2} compute the exact solution of
$H\Delta=g$ for the SPD block-cyclic-tridiagonal $H$ of \eqref{eq:HBT}, using
\[
  O(P\,m^3)\ \text{flops and}\ O(P\,m^2)\ \text{storage},
\]
i.e.\ \emph{linear} in the period $P$.
\end{lemma}

\begin{proof}
\emph{Correctness.} By construction $H=T+UKU^{\!\top}$ with $U,K$ as in
\eqref{eq:UKV}. The block recursion \eqref{eq:thomas} is the exact block LU/LDL
solve of the block-tridiagonal system $T\Delta=g$: substituting the
back-substitution formula $\Delta_k=M_k^{-1}(y_k-E_k\Delta_{k+1})$ and the
definitions of $M_k,y_k$ into the $k$-th block row
$E_{k-1}^{\!\top}\Delta_{k-1}+\widetilde D_k\Delta_k+E_k\Delta_{k+1}=g_k$ verifies the
identity for every $k$. Hence $T^{-1}$ is computed exactly, and the
Woodbury identity \eqref{eq:woodbury} (or \eqref{eq:woodbury2} when $K$ is singular)
is an algebraic identity valid for any invertible $H,T$ with $H=T+UKU^{\!\top}$.
Since $H\succ0$ on $\mathcal F$ by Lemma~\ref{lem:struct} and $T$ (or $T_\rho$) is
invertible by the pivot/regularisation construction, \eqref{eq:woodbury2} returns the
exact $H^{-1}g$. (Numerical exactness up to roundoff is verified in
Remark~\ref{rem:num}.)

\emph{Complexity.} The forward sweep performs, per index $k$, one $m\times m$
symmetric factorisation and a constant number of $m\times m$ matrix products:
$O(m^3)$; over $k=0,\dots,P-1$ this is $O(Pm^3)$. Applying $T^{-1}$ to $g$ and to the
$2m$ columns of $U$ costs $O(Pm^3)$ once the factors $\{M_k\}$ are stored (each apply
is $O(Pm^2)$, and there are $1+2m$ right-hand sides). The $2m\times2m$ inner solve is
$O(m^3)$. Storage holds the $P$ factored blocks $M_k$ ($O(Pm^2)$) and the
right-hand sides; no $mP\times mP$ object is ever formed.
\end{proof}

\begin{remark}[Numerical verification]\label{rem:num}
We implemented \eqref{eq:thomas}--\eqref{eq:woodbury} for random SPD
block-cyclic-tridiagonal systems ($P=6$, $m=3$) and compared against a dense direct
solve; the relative error was $5.1\times10^{-16}$, i.e.\ machine precision,
confirming exactness and stability of the structured solve.
\end{remark}

\section{The algorithm}\label{sec:algo}

We combine the barrier of \S\ref{sec:barrier} with the structured solver of
\S\ref{sec:linalg} in a short-step path-following method, preceded by a Phase-I that
either produces a strictly feasible starting point or certifies infeasibility.

\subsection{Phase~I: feasibility decision and a feasible start}\label{sec:phaseI}

Introduce a scalar slack $\tau$ and consider the auxiliary periodic LMI
\begin{equation}\label{eq:phaseI}
  \mathcal A_k(\bm x)+\tau I_n\ \succ\ 0,\qquad k=0,\dots,P-1,
\end{equation}
together with the bounding box $\norm{\bm x}^2+\tau^2\le R^2$ of Assumption~(A).
For $\tau$ large enough \eqref{eq:phaseI} is strictly feasible (take $\bm x=0$,
$\tau>\max_k\norm{F_0(k)}$), giving a trivial strictly feasible start for the
\emph{auxiliary} barrier
\[
  F_\tau(\bm x,\tau)=-\sum_{k=0}^{P-1}\log\det\big(\mathcal A_k(\bm x)+\tau I_n\big)
                     -\log\big(R^2-\norm{\bm x}^2-\tau^2\big).
\]
The augmented residual $\mathcal A_k(\bm x)+\tau I_n$ still depends only on
$\big(x(k),x(k+1),\tau\big)$, so Lemma~\ref{lem:struct} applies verbatim with one
extra scalar variable $\tau$ that couples to \emph{all} blocks; $\tau$ is therefore a
single bordering row/column and is eliminated by one more rank-$1$
Sherman--Morrison step appended to \eqref{eq:woodbury}, preserving the $O(Pm^3)$
cost. Minimising $\tau$ subject to \eqref{eq:phaseI} along the central path of
$F_\tau$ yields the optimal slack
\[
  \tau^\star\;:=\;\inf\{\tau:\ \exists\,\bm x\ \text{with }\mathcal A_k(\bm x)+\tau I\succ0\ \forall k,\ \norm{\bm x}\le R\}.
\]
Then $\mathcal F\cap\{\norm{\bm x}\le R\}\neq\emptyset$ \emph{iff} $\tau^\star\le0$,
with $\tau^\star<0$ giving a strictly interior point; $\tau^\star>0$ certifies
infeasibility within the box. This decides (i). The minimiser at any $\tau<0$ is a
strictly feasible $\bm x$, providing the start for Phase~II and answering (ii).

\subsection{Phase~II: path following}

\begin{algorithm}[t]
\caption{Periodic-structure interior-point solver for PLMI}
\label{alg:main}
\begin{algorithmic}[1]
\Require periodic data $\{F_i(k),\,C(k,\cdot,\cdot)\}_{k=0}^{P-1}$; tolerance
         $\varepsilon>0$; box radius $R$ (Assumption~(A)); barrier parameter
         $\nu\le 2nP$ from \eqref{eq:nu}.
\Ensure either ``infeasible'' with certificate $\tau^\star>0$, or a strictly
         feasible $\bm x=(x(0),\dots,x(P-1))$, equivalently $X(0),\dots,X(P-1)$.
\State \textbf{Phase I.} Run the path-following loop below on $F_\tau$ from the start
       $(\bm x,\tau)=(0,\;1+\max_k\norm{F_0(k)})$ until $\tau$ crosses $0$ or the
       central-path limit certifies $\tau^\star\ge0$.
\If{$\tau^\star>0$} \State \Return ``infeasible'', certificate $\tau^\star$.
\Else \State set $\bm x^{(0)}\gets$ current strictly feasible point; $t\gets t_0>0$.
\EndIf
\State \textbf{Phase II.}
\While{$\nu/t>\varepsilon$}
  \State assemble blocks $D_k,E_k$ of $H=\nabla^2F(\bm x)$ and $g=\nabla F(\bm x)$
         via \eqref{eq:grad}--\eqref{eq:hess}, using only $S_k=\mathcal A_k(\bm x)$ and
         its Cholesky factor; \Comment{$O(P(m^2n^2+mn^3+n^3))$ flops}
  \State solve $H\,\Delta=-g$ by \eqref{eq:thomas}--\eqref{eq:woodbury2};
         \Comment{$O(Pm^3)$ flops (Lemma~\ref{lem:exact})}
  \State Newton decrement $\lambda\gets\sqrt{-g^{\!\top}\Delta}$
  \If{$\lambda\le \tfrac14$} \State $\bm x\gets\bm x+\Delta$; $t\gets(1+\tfrac{1}{8\sqrt{\nu}})\,t$
  \Else \State $\bm x\gets\bm x+\tfrac{1}{1+\lambda}\Delta$ \Comment{damped Newton step}
  \EndIf
\EndWhile
\State \Return $\bm x$.
\end{algorithmic}
\end{algorithm}

For a pure feasibility problem the cost vector $c$ in the centering objective
$t\,c^{\!\top}\bm x+F(\bm x)$ is zero, so Phase~II simply tracks the analytic centre of
$\mathcal F$; the increment of $t$ is harmless but kept so the algorithm also solves
the optimisation variant $\min c^{\!\top}\bm x$ over $\mathcal F$ verbatim.

\section{Convergence and complexity}\label{sec:conv}

\begin{theorem}[Global convergence and iteration bound]\label{thm:conv}
Under Assumption~(A), Algorithm~\ref{alg:main} terminates. Phase~I decides
feasibility correctly: it returns ``infeasible'' iff $\mathcal F\cap\{\norm{\bm
x}\le R\}=\emptyset$. When feasible, Phase~II returns a strictly feasible point, and
the total number of Newton steps to reach duality gap $\le\varepsilon$ (for the
optimisation variant) or to centre to accuracy $\varepsilon$ (for feasibility) is
\begin{equation}\label{eq:itbound}
  N_{\mathrm{Newton}}
  \;=\;O\!\Big(\sqrt{\nu}\,\log\tfrac{\nu}{t_0\,\varepsilon}\Big)
  \;=\;O\!\Big(\sqrt{nP}\,\log\tfrac{1}{\varepsilon}\Big).
\end{equation}
\end{theorem}

\begin{proof}
By Lemma~\ref{lem:sc} the barrier $F$ (and $F_\tau$ in Phase~I) is self-concordant
with parameter $\nu\le2nP$. Algorithm~\ref{alg:main} is the textbook short-step
barrier method (Nesterov--Nemirovski; or Boyd--Vandenberghe, \emph{Convex
Optimization}, \S11.5; Renegar, \emph{A Mathematical View of Interior-Point
Methods}, Thm.~2.4.1) applied to that barrier. Two standard facts:

(1) \emph{Centering preserves the near-central condition.} If the current iterate
satisfies the proximity bound $\lambda(\bm x;t)\le\frac14$ (Newton decrement of
$t\,c^{\!\top}\bm x+F$), one full Newton step gives $\lambda^+\le2\lambda^2\le
2(\tfrac14)^2=\tfrac18\le\frac14$ (the self-concordant quadratic-convergence
estimate, Nesterov--Nemirovski Thm.~2.2.3). Thus the loop maintains
$\lambda\le\frac14$ after the first centring.

(2) \emph{The $t$-update stays in the region of quadratic convergence.}
Increasing $t$ by the factor $1+\frac{1}{8\sqrt\nu}$ raises the Newton decrement of
the new $t$ to at most a constant below $\frac14$ (Nesterov--Nemirovski Prop.~5.1.4
/ Renegar Thm.~2.4.1), so a single centring step restores $\lambda\le\frac14$. Hence
each outer iteration is one Newton step.

Because $t$ grows geometrically with ratio $1+\frac1{8\sqrt\nu}$, the number of outer
iterations to increase $t$ from $t_0$ to $\nu/\varepsilon$ is
\[
  \frac{\log(\nu/(t_0\varepsilon))}{\log(1+\tfrac1{8\sqrt\nu})}
  = O\!\big(\sqrt\nu\,\log\tfrac{\nu}{t_0\varepsilon}\big),
\]
which is \eqref{eq:itbound} since $\nu\le2nP$. The duality gap on the central path is
exactly $\nu/t$ (Boyd--Vandenberghe \S11.2.2), so reaching $t\ge\nu/\varepsilon$
guarantees gap $\le\varepsilon$; for feasibility, $\lambda\le\frac14$ certifies a
strictly feasible point. The damped-Newton fallback ($s=\frac1{1+\lambda}$) gives the
standard guaranteed decrease $F(\bm x^+)-F(\bm x)\le-(\lambda-\log(1+\lambda))<0$ for
the at most $O(1)$ initial iterates where $\lambda>\frac14$, so the method reaches the
quadratically convergent region in finitely many steps (Boyd--Vandenberghe \S9.6.4).
Phase~I correctness follows from the definition of $\tau^\star$ in \S\ref{sec:phaseI}:
the auxiliary problem is strictly feasible for large $\tau$, its optimal value
$\tau^\star$ exists by compactness of the box, and $\tau^\star\le0\Leftrightarrow$ a
feasible $\bm x$ exists in the box, $\tau^\star>0\Leftrightarrow$ no feasible $\bm x$
exists in the box. All hypotheses of the cited theorems (self-concordance, nonempty
bounded interior) hold by Lemma~\ref{lem:sc} and Assumption~(A).
\end{proof}

\begin{theorem}[Per-iteration complexity and total cost]\label{thm:cost}
Each Newton step of Algorithm~\ref{alg:main} costs
\begin{equation}\label{eq:periter}
  O\!\big(P\,(m^2 n^2 + m\,n^3 + n^3)\big)\ \text{flops}
  \qquad\text{and}\qquad
  O\!\big(P\,(m+n^2)\big)\ \text{storage},
\end{equation}
both \emph{linear} in the period $P$. Combined with Theorem~\ref{thm:conv}, the total
work to solve the PLMI to accuracy $\varepsilon$ is
\begin{equation}\label{eq:total}
  O\!\Big(\sqrt{nP}\;P\,(m^2n^2+m\,n^3+n^3)\,\log\tfrac1\varepsilon\Big)
  = O\!\Big(P^{3/2}\,n^{1/2}\,(m^2n^2+mn^3+n^3)\,\log\tfrac1\varepsilon\Big).
\end{equation}
In particular the dependence on $P$ is $P^{3/2}$, versus $P^{3}$ (or worse) for the
naive lifted dense solver, and storage is $O(P)$ versus $O(P^2)$.
\end{theorem}

\begin{proof}
\emph{Assembly \eqref{eq:grad}--\eqref{eq:hess}.} For each $k$ we factor
$S_k\in\Sym^n$ once: Cholesky in $O(n^3)$; total $O(Pn^3)$. The derivative matrices
$\partial_{x_i(k)}S_\ell$ are the fixed data, already available. For the Hessian
block coupling instant $\ell$, precompute the $m$ matrices
$W_i^\ell:=S_\ell^{-1}\partial_{x_i(\ell)}S_\ell$ and
$W_i^{\ell,+}:=S_\ell^{-1}\partial_{x_i(\ell+1)}S_\ell$ by $O(m)$ triangular solves
with $n$ right-hand sides each, costing $O(mn^3)$ per $\ell$. Then each Hessian
entry \eqref{eq:hess} restricted to instant $\ell$ is a trace
$\tr(W_i^\ell W_j^\ell)$ (or mixed) costing $O(n^2)$, and there are $O(m^2)$ such
entries in the $O(1)$ blocks touching $\ell$: $O(m^2n^2)$ per $\ell$. The gradient
\eqref{eq:grad} costs $O(mn^2)$ per $\ell$, dominated. Summing over
$\ell=0,\dots,P-1$ gives $O\!\big(P(mn^3+m^2n^2)\big)$ plus the $O(Pn^3)$ Cholesky
term, i.e.\ the flop bound \eqref{eq:periter}. Storage: the $P$ Cholesky factors and
Hessian blocks are $O(P(n^2+m^2))$; writing $O(P(m+n^2))$ covers the natural
matrix-variable regime $m\le\binom{n+1}{2}=O(n^2)$ as well as small $m$.

\emph{Newton solve.} By Lemma~\ref{lem:exact}, $O(Pm^3)$ flops, $O(Pm^2)$ storage,
dominated by the assembly cost whenever $m\lesssim n^2$ (otherwise add the linear-in-$P$
term $O(Pm^3)$).

\emph{Total.} Multiplying \eqref{eq:periter} by the iteration count \eqref{eq:itbound}
with $\sqrt\nu=O(\sqrt{nP})$ gives \eqref{eq:total}. Comparison with the lift: a
dense solve of the $mP\times mP$ Hessian is $O((mP)^3)=O(m^3P^3)$ per step with
$O(m^2P^2)$ storage; our method replaces $P^3\to P^{3/2}$ in the $P$-dependence and
$P^2\to P$ in storage.
\end{proof}

\section{Structure-preserving (periodic Schur / periodic QZ) realisation and stability}
\label{sec:schur}

Requirement (iii) asks that the block solves admit a structure-preserving,
numerically stable implementation. Two points complete this.

\paragraph{(a) Periodic Schur for residual handling and warm starts.}
In the Lyapunov / KYP instances the residuals are
$S_k=-(A(k)^{\!\top}X(k+1)A(k)-X(k)+Q(k))$. Forming and refactoring $S_k$ at each
Newton step naively recomputes products $A(k)^{\!\top}X(k+1)A(k)$. Instead, compute
once the \emph{periodic real Schur decomposition} of the matrix sequence
$\{A(k)\}_{k=0}^{P-1}$: orthogonal $\{U_k\}$ with $U_{k+P}=U_k$ such that
$\widehat A(k):=U_{k+1}^{\!\top}A(k)U_k$ is (quasi-)upper-triangular for all $k$ and
the product $\widehat A(P-1)\cdots\widehat A(0)$ is the real Schur form of the
monodromy matrix $\Phi=A(P-1)\cdots A(0)$. This decomposition is computed by the
periodic QR/QZ algorithm in $O(Pn^3)$ flops and is \emph{backward stable}
(orthogonal transformations only); see Bojanczyk--Golub--Van~Dooren (1992) and
Kressner (2006). Working in the rotated coordinates $\widehat X(k):=U_k^{\!\top}
X(k)U_k$ turns the dense coupling $A(k)^{\!\top}X(k+1)A(k)$ into the triangular
product $\widehat A(k)^{\!\top}\widehat X(k+1)\widehat A(k)$, so each residual update
is triangular and costs $O(n^3)$ with no growth in $P$, while all positive
definiteness tests are invariant under the orthogonal change of basis. This realises
(iii) and keeps the per-iteration cost \eqref{eq:periter}.

\paragraph{(b) Backward stability of the structured Newton solve.}
The block sweep \eqref{eq:thomas} is the block LDL$^{\!\top}$ factorisation of the SPD
matrix $T$; carried out with symmetric (Bunch--Kaufman) pivoting inside each
$m\times m$ block factorisation, it is backward stable with a growth factor bounded
independently of $P$ (each pivot $M_k$ is a Schur complement of an SPD matrix, hence
positive definite and well conditioned relative to $H$). The Woodbury correction
\eqref{eq:woodbury}--\eqref{eq:woodbury2} adds only a $2m\times2m$ (or, with the
$\tau$ border, $(2m{+}1)\times(2m{+}1)$) solve; its conditioning is governed by
$K^{-1}+U^{\!\top}T^{-1}U$, a principal compression of $H^{-1}$, which therefore
inherits the conditioning of $H=\nabla^2F\succ0$. Because the dominant computation
uses only orthogonal transformations (periodic Schur) and SPD block-Cholesky/LDL
with pivoting, the computed step $\widehat\Delta$ solves
$(H+\delta H)\widehat\Delta=g+\delta g$ with $\norm{\delta H}=O(u)\norm H$,
$\norm{\delta g}=O(u)\norm g$ and constants independent of $P$, where $u$ is the unit
roundoff. This is the provable numerical stability requested.

\section{Summary of what is established}\label{sec:summary}

\begin{itemize}
\item \textbf{(i) Feasibility decision:} Phase~I (\S\ref{sec:phaseI}) computes the
optimal slack $\tau^\star$; $\mathcal F\cap\{\norm{\bm x}\le R\}\neq\emptyset$ iff
$\tau^\star\le0$, with $\tau^\star>0$ a verified infeasibility certificate
(Theorem~\ref{thm:conv}).
\item \textbf{(ii) Periodic solution:} when feasible, the centred Phase~II iterate is
a strictly feasible $P$-periodic $X(0),\dots,X(P-1)$ (Theorem~\ref{thm:conv}).
\item \textbf{(iii) No $nP\times nP$ lift; structure preserving:} the only large
object touched is the SPD block-cyclic-tridiagonal Hessian, solved in $O(Pm^3)$ by
\eqref{eq:thomas}--\eqref{eq:woodbury2} (Lemma~\ref{lem:exact}); residual algebra is
done in periodic Schur coordinates (\S\ref{sec:schur}).
\item \textbf{Convergence:} $O(\sqrt{nP}\log\frac1\varepsilon)$ Newton steps
(Theorem~\ref{thm:conv}).
\item \textbf{Per-iteration complexity:} $O\!\big(P(m^2n^2+mn^3+n^3)\big)$ flops,
$O(P(m+n^2))$ storage, both linear in $P$; total
$O\!\big(P^{3/2}n^{1/2}(m^2n^2+mn^3+n^3)\log\frac1\varepsilon\big)$
(Theorem~\ref{thm:cost}), strictly better in $P$ than the lifted $O(m^3P^3)$ solver.
\item \textbf{Stability:} orthogonal periodic Schur transformations plus pivoted SPD
block factorisation give backward stability with $P$-independent constants
(\S\ref{sec:schur}).
\end{itemize}

\end{document}
